%%%%%%%%%%%%%%%%%%%%%%%%%%%%%%%%%%%%%%%%%
% "ModernCV" CV and Cover Letter
% LaTeX Template
% Version 1.1 (9/12/12)
%
% This template has been downloaded from:
% http://www.LaTeXTemplates.com
%
% Original author:
% Xavier Danaux (xdanaux@gmail.com)
%
% License:
% CC BY-NC-SA 3.0 (http://creativecommons.org/licenses/by-nc-sa/3.0/)
%
% Important note:
% This template requires the moderncv.cls and .sty files to be in the same 
% directory as this .tex file. These files provide the resume style and themes 
% used for structuring the document.
%
%%%%%%%%%%%%%%%%%%%%%%%%%%%%%%%%%%%%%%%%%

%----------------------------------------------------------------------------------------
%	PACKAGES AND OTHER DOCUMENT CONFIGURATIONS
%----------------------------------------------------------------------------------------

\documentclass[10pt,a4paper]{moderncv} % Font sizes: 10, 11, or 12; paper sizes: a4paper, letterpaper, a5paper, legalpaper, executivepaper or landscape; font families: sans or roman

\moderncvstyle{classic} % CV theme - options include: 'casual' (default), 'classic', 'oldstyle' and 'banking'
\moderncvcolor{grey} % CV color - options include: 'blue' (default), 'orange', 'green', 'red', 'purple', 'grey' and 'black'
\usepackage{lipsum}
\usepackage{amsmath}
\usepackage{url}
\usepackage[scaled]{beramono}
\renewcommand*\familydefault{\ttdefault} %% Only if the base font of the document is to be typewriter style
\usepackage[T1]{fontenc}

\usepackage[scale=0.87]{geometry} % Reduce document margins
\linespread{1.1}
\setlength{\hintscolumnwidth}{2.16cm} % Uncomment to change the width of the dates column
%\setlength{\makecvtitlenamewidth}{10cm} % For the 'classic' style, uncomment to adjust the width of the space allocated to your name

%----------------------------------------------------------------------------------------
%	NAME AND CONTACT INFORMATION SECTION
%----------------------------------------------------------------------------------------

\firstname{K{\i}van\c{c}} % Your first name
\familyname{Tatar} % Your last name

% All information in this block is optional, comment out any lines you don't need
\title{Assistant Professor in Interactive AI}
\address{Chalmers University of Technology \\Data Science and AI Division \\ Computer Science Department \\ Rännvägen 6B,\\}{  
412 58 Gothenburg, Sweden}

\email{tatar@chalmers.se}
\homepage{aicomparts.com}{Research Group} 
\webpage{kivanctatar.com}{Personal Website} 

%----------------------------------------------------------------------------------------
\begin{document}
	%\ttfamily
	%\asciifamily
	\makecvtitle % Print the CV title
\begin{center} \LARGE  Description of Time Distribution\end{center}
\hline

\cvitem{}{My contract is on 80\% research and 20\% teaching. My research time is funded by WASP-HS.}
\cvitem{}{\textbf{I have written, prepared, and submitted my Docent and Associate Professor application} to the CSE department, which is now approved by the department to be sent to the appointment committee at Chalmers central. Writing my docent application and the pedagogical portfolio took significant amount of my time in 2025.}

\begin{center} \LARGE  Research Experience and Merits\end{center}
\hline

    \section{Research Grants}

    \cvitem{Summary}{I have acquired two significant grants as co-PI, totalling to a 48 million SEK. It is also important to note that I am the only WASP-HS assistant prof. who has acquired a WASP-HS cluster grant as a co-PI or PI.}
    
    {\begin{center}{\textcolor{magenta}{Submitted Grants Applications Still Awaiting Decisions}}\end{center}}

    \cvitem{VR-2}{\href{https://vr.se/en/}{Vetenskapsrådet}, \href{https://www.vr.se/english/applying-for-funding/calls/2024-11-12-research-project-grant-within-natural-and-engineering-sciences.html}{\textbf{Research Project Grant in Natural Sciences and Engineering}}. \textit{Resonant Spaces: Mathematical Foundations of Deep Generative Modelling for Audio Signals}. Kıvanç Tatar (PI, Chalmers University of Technology and University of Gothenburg), Kathlén Kohn (co-PI, KTH Royal Institute of Technology), Ashkan Panahi (co-PI, Chalmers University of Technology and University of Gothenburg). 5.7 million SEK.}

    
    \cvitem{}{\begin{center}\textcolor{magenta}{\textcolor{magenta}{Successful Grants}}\end{center}}
    \cvitem{WASPHS-3}{\href{https://wasp-hs.org/open-call-for-wasp-hs-research-clusters/}{WASP-HS Research Cluster Call}. \textit{AI Futures of Culture and Memory}. Anna Foka (PI, Uppsala University), Andre Holzapfel (co-PI, KTH Royal Institute of Technology), Kıvanç Tatar (co-PI, Chalmers University of Technology and University of Gothenburg), Coppelie Cocq (co-PI, Umeå University) Karolja Golub (co-PI, Linnaeus University). \textbf{30 million SEK}.}
    
    \cvitem{VR-1}{\href{https://vr.se/en/}{Vetenskapsrådet}, \href{https://www.vr.se/english/applying-for-funding/calls/2023-11-15-research-environment-grant---humanities-and-social-sciences.html}{Research environment grant – humanities and social sciences}. \textit{VOICE. AI-generated voices. Legal and societal perspectives.} Katja de Vries (PI, Uppsala University), Andre Holzapfel (co-PI, KTH Royal Institute of Technology), Kıvanç Tatar (co-PI, Chalmers University of Technology and University of Gothenburg), Kalle Berglund (co-PI, Uppsala University). \textbf{18 million SEK}.}

    \cvitem{}{\begin{center}\textcolor{magenta}{\textcolor{magenta}{Rejected Applications}}\end{center}}

    \cvitem{SSF-1}{SSF Future Leaders Grant. 18 million SEK. This took me four weeks to write (whole August 2024), and unfortunately, my application was rejected.}

\section{Academic Publications}

\cvitem{Summary}{\textbf{I have published two journal papers, and an additional journal paper is accepted and in press.} Additionally, three journal papers are submitted. It is exceptional that I have submitted 5 journal papers in 2025.}
    
	\cvitem{}{\begin{center} \textcolor{magenta}{\textcolor{magenta}{Journal Papers}}\end{center}}

    \cvitem{1}{Zappi, V.,\textcolor{magenta}{Tatar, K}. (2025). \textit{Neural audio instruments: Epistemological and phenomenological perspectives on musical embodiment of deep learning}. \textbf{Frontiers in Computer Science}, 7. \url{https://doi.org/10.3389/fcomp.2025.1575168}}

    \cvitem{2}{Cotton, K., Kaila, A.K.,Jääskeläinen, P., Holzapel, A., \textcolor{magenta}{Tatar, K}. (2025).\textit{Imploding between the Facts and Concerns of Musical-AI: A Multidisciplinary Method for Analysing Human-AI Musical Interaction.} \textbf{Humanities and Social Sciences Communications}, 12(1), 1–20, \textbf{Springer Nature Press}. \url{https://doi.org/10.1057/s41599-025-04533-4}}

    \cvitem{3}{Madaghiele V., Lund L., Holzer D., Kelkar T., \textcolor{magenta}{Tatar, K.}, and Holzapfel A. \textit{Expanding the machine: notating generative synthesis with a state-based representation and an interactive timbre space}. \textbf{Accepted for publications and in press} at Organised Sound, Cambridge Press.}

    \cvitem{4}{Cotton, K., de Vries K., Holzapfel A., Berglund K., \textcolor{magenta}{Tatar, K.}. \textit{Revising MIR Research Practices for Singing Data Collection}. Submitted to AI and Society, Springer Nature.}

    \cvitem{5}{Liu X., Cotton K., \textcolor{magenta}{Tatar, K.} \textit{Revisiting a Decade of Computational Game Creativity}. Submitted to International Journal on Human Computer Interaction, Taylor and Francis. Passed the editor and now in review.}

    \cvitem{6}{*Fröhlich M. M., Cotton K., \textcolor{magenta}{Tatar, K.} \textit{Artificial Intelligence: Perspectives on a Caesura in Art, its Creation and Consumption}. Submitted to Leonardo, MIT Press.}    

	\cvitem{}{\small All papers are double-blind peer-reviewed unless otherwise is indicated.}
	\cvitem{}{\small *These papers are single-blind peer-reviewed. The reviewers had access to the list of authors.}

\begin{center} \LARGE  Teaching Experience and Merits\end{center}
\hline
    
\section{Supervision Experience}

\cvitem{Summary}{Kelsey Cotton has passed licentiate defense.}

\cvitem{}{\begin{center}\textcolor{magenta}{\textcolor{magenta}{Post-doctoral Fellows}}\end{center}}

\cvitem{Ongoing}{\href{https://scholar.google.com/citations?user=ni5MM4oAAAAJ}{\textbf{Xuechen (Hugh) Liu}}. \textit{Postdoc in Interactive AI for interdisciplinary Artistic Practices}. 2023-01::2026-01. \textcolor{magenta}{Supervisor: K{\i}van\c{c} Tatar}. }
\cvitem{}{\begin{center}\textcolor{magenta}{\textcolor{magenta}{Ph.D. students}}\end{center}}

\cvitem{Ongoing}{\href{https://kelseycotton.com/}{\textbf{Kelsey Cotton}}. \textit{Interactive Music Systems using Artificial Intelligence}. Expected graduation 2027. \textcolor{magenta}{Main Supervisor: K{\i}van\c{c} Tatar}, Co-supervisor: Devdatt Dubhashi, Examiner: Graham Kemp. Funded by WASP-HS-1. Licentiate defense is scheduled to 2024-11-26 with the external examiner Prof. Alexander Refsum Jensenius from University of Oslo, Norway.}

\cvitem{}{\textbf{New PhD}. \textit{Generative AI in Culture and Creative Domains}. PhD hiring is ongoing. We are now interviewing shortlisted candidates. \textcolor{magenta}{Main Supervisor: K{\i}van\c{c} Tatar}, Co-supervisor: Ashkan Panahi, Examiner: Dag Wedelin.}

\cvitem{}{\begin{center}\textcolor{magenta}{\textcolor{magenta}{Master students}}\end{center}}

\cvitem{Summary}{I have supervised two master thesis in 2025. These two theses are finalized now. I have examined 3 master theses in addition. Two are finalized now.}

\cvitem{2025}{Qi Chen and Gritta Joshy. \textit{Curriculum Learning for Symbolic Music Generation with Transformers}. \textcolor{magenta}{Main Supervisor: K{\i}van\c{c} Tatar}, co-Supervisor: Stefano Sarao Mannelli, examiner: Dag Wedelin. Master Thesis.}

\cvitem{2025}{Ipek Korkmaz. \textit{Investigating Impact of Audio Features Loss Functions for Audio Synthesis using Variational Autoencoders}. \textcolor{magenta}{Main Supervisor: K{\i}van\c{c} Tatar}, examiner: Simon Olsson. Master Thesis.}

\section{Teaching Experience}

\cvitem{2025 Spring}{\textbf{Examiner.} Advanced course. TRA 385 - Machine Learning and AI through artistic innovation, Chalmers University of Technology, Gothenburg, Sweden. This course is produced by K{\i}van\c{c} Tatar, with contributions from Kelsey Cotton and Hugh Liu Xuechen, as a part of \href{https://www.chalmers.se/en/education/your-studies/course-selection-and-registration/select-courses/choose-a-tracks-course/}{TRACKS} initiative, conducted using the \href{https://www.chalmers.se/en/education/study-at-chalmers/why-chalmers/fuse-chalmers-makerspace/}{FUSE} maker spaces of Chalmers. }

\cvitem{}{In 2025, TRA385 Tracks course is updated to Machine Learning and AI through Artistic Innovations, inline with the comments from the CSE IB committee at my half-time evaluation. The course received quite postive feedback from the students in 2025 run.}

\cvitem{New Course}{Masters in Acoustics program approaches me to initiate a new course in Musical AI and Music Information Retrieval. I have been in communication with DSAI teaching management (Birgit Grohe and Ashkan Panahi), who have been positive to this course. Ideally, this course will be an elective for both Masters in DSAI and Masters in Acoustics programs, thus reaching to a minimum of two masters programs. The course can run in 2028, so there is still time to plan.}

\section{Personal Development}

\cvitem{}{I have been learning a variety of workshop skills towards projects in physical computing. Last year, I have took an introduction to metal workshop. Additionally, I started building modular synths using pre-made packages. In future years, I am planning to design my own audio synth circuits to come up with future research projects in audio circuit design using machine learning.}

\begin{center} \LARGE Academic Citizenship\end{center}
\hline

\section{Reviewing for Funding Agencies and Other Universities}
\cvitem{2025}{\textbf{External Examiner}. Reviewed one master thesis for the masters in \href{https://www.uio.no/english/studies/programmes/mct-master/}{Music, Communication and Technology} at University of Oslo in Norway.}
\cvitem{2025}{\textbf{Reviewer}. Reviewed a grant application for the \textbf{Swiss National Science Foundation}.}
\cvitem{2024}{\textbf{Reviewer}. Reviewed a grant application for the \textbf{KU Leuven Research Council} in Belgium.}
\cvitem{2022}{\textbf{External Examiner}. Reviewed one master thesis for the masters in \href{https://www.uio.no/english/studies/programmes/mct-master/}{Music, Communication and Technology} at University of Oslo in Norway.}
\cvitem{2022}{\textbf{Jury duty}. Reviewed 96 grant applications for Concept to Realization Track at the Canada Council for the Arts.}

\section{Activities at Chalmers}
\cvitem{1}{New Compute Cluster for teaching at Data Science and AI division, pursued initially by Matti Karppa, Emilio Jorge, and K{\i}van\c{c} Tatar. This initiative is now successful, and the CSE department with additional budget from Chalmers centrally decided to purchase a new computing cluster for DSAI teaching thanks to this initiative. Hardware purchase and planning is ongoing as of 2024-09, with expected completion before 2025.}

\cvitem{2}{I have led an initiative to update the DSAI division page on Chalmers website so that the web content represents the depth, breath, and diversity of research at DSAI. The faculty discussions are done as of 2024-09, updates are now being added to the website. }

\cvitem{3}{I contributed majorly to the FUSE learning environments. With my inputs and planning of TRA385, the main FUSE area is set up to be a collaborative and open studio-based learning environment, where students from different course have knowledge exchanges. This setup was a direct result of my engagement with FUSE staff and management. They were so inspired by the idea that this is now a status-quo learning environment setup for the FUSE main hall.}

\cvitem{4}{I contributed to building FUSE electronics lab from scratch with its requirements, devices, and components. The electronics lab is still improving with my inputs.}

\cvitem{5}{I have joined the TRACKs Teachers' Day in Fall 2024, a full day event covering activities and pedagogy at TRACKs.}

\cvitem{6}{I run an example workshop from the TRA385 course for the Redir event (“Rektorer och dekaner i riket”) at Chalmers, which is a formal event where vice chancellors and people in high positions from universities from all over Sweden came to Chalmers. The event was on May 16th, 2025.}

\section{Reviewing Activities for Journals and Conferences}

\cvitem{2025}{\textbf{Reviewer}. \href{https://www.frontiersin.org/journals/computer-science}{Discover Artificial Intelligence, Springer Nature Press}. 1 paper}
\cvitem{2025}{\textbf{Reviewer}. \href{https://link.springer.com/journal/44163}{Frontiers in Computer Science, Frontiers Press}. 1 paper}
\cvitem{2025}{\textbf{Meta-Reviewer}. \href{https://nime2025.org/}{New-Interfaces for Musical Expression Conference 2025}. 4 papers}
\cvitem{2025}{\textbf{Reviewer}. \href{https://link.springer.com/journal/521}{Neural Computing and Applications}, Springer Press, 1 paper.}
\cvitem{2024}{\textbf{Reviewer}. \href{https://direct.mit.edu/leon}{Leonardo}, MIT Press, 1 paper.}

\section{Other}

\cvitem{}{I have joined four ERC information seminars and workshops covering the ERC Starting grant.}


\begin{center} \LARGE Research Utilization\end{center}
\hline

\section{Artworks exhibited in Public Venues}

\cvitem{2024-08-31}{\href{https://kivanctatar.com/Exposing-the-Bias-in-AI}{\textbf{Exposing the Bias in Artificial Intelligence}}. An audiovisual live performance at the \href{https://koami.art/}{Koami Arts Festival}, 300 in-person audience, Gothenburg, Sweden.}


\section{Press Coverage}

\cvitem{2025-05-18}{I gave an interview to \textbf{Sveriges Radio P3 Nyheter} on AI-slop. Excerpts from the interview were shared on the radio, available at\url{https://www.sverigesradio.se/artikel/ar-ai-slop-slutet-for-sociala-medier}.}



\end{document}